\documentclass[11pt,a4paper,fontset=none]{ctexart}
\usepackage[margin=2.3cm]{geometry}
\usepackage{amsmath,amssymb,mathtools}
\usepackage{graphicx}
\usepackage{booktabs}
\usepackage[table]{xcolor}
\usepackage{hyperref}
\usepackage[most]{tcolorbox}
\usepackage[numbers,sort&compress]{natbib}
\usepackage{array}

\definecolor{accentblue}{HTML}{2F5D80}
\definecolor{softblue}{HTML}{EEF5FA}
\definecolor{linegray}{HTML}{D8E3EC}
\definecolor{darktext}{HTML}{243746}

\hypersetup{
  colorlinks=true,
  linkcolor=accentblue,
  citecolor=accentblue,
  urlcolor=accentblue,
  hypertexnames=false,
  pdftitle={MultiSpec-GeoDiff: 多模态谱图驱动的分子结构反演 Demo},
  pdfauthor={OpenAI Codex for DeepTech Demo Proposal}
}

\setlength{\parindent}{2em}
\setlength{\parskip}{0.35em}
\linespread{1.12}

\setmainfont{TeX Gyre Termes}
\setsansfont{TeX Gyre Heros}
\setmonofont{Latin Modern Mono}
\setCJKmainfont[AutoFakeBold=2,AutoFakeSlant=0.2]{Microsoft YaHei}
\setCJKsansfont[AutoFakeBold=2]{Microsoft YaHei}
\setCJKmonofont{DengXian}
\setCJKfamilyfont{zhkai}{KaiTi}
\newcommand{\kai}{\CJKfamily{zhkai}}

\newtcolorbox{keybox}[1][]{
  enhanced,
  breakable,
  colback=softblue,
  colframe=accentblue,
  boxrule=0.8pt,
  arc=1.5mm,
  left=1.5mm,
  right=1.5mm,
  top=1.0mm,
  bottom=1.0mm,
  title=#1,
}

\newcommand{\placeholderfigure}[1]{%
  \begin{tcolorbox}[
    width=0.92\linewidth,
    height=5.2cm,
    colback=white,
    colframe=accentblue!80!black,
    boxrule=0.8pt,
    arc=1.2mm,
    valign=center,
  ]
  \centering
  {\large\bfseries\color{accentblue}#1}
  \end{tcolorbox}%
}

\newcommand{\projectname}{MultiSpec-GeoDiff}
\newcommand{\projecttitle}{\projectname: 多模态谱图驱动的分子结构反演 Demo}

\begin{document}

\begin{titlepage}
  \centering
  \vspace*{2.8cm}
  {\Huge\bfseries\color{accentblue}\projectname\par}
  \vspace{0.4cm}
  {\LARGE\bfseries 多模态谱图驱动的分子结构反演 Demo\par}
  \vspace{1.0cm}
  {\large DP Technology / Deep Potential 实习申请 Demo Proposal Package\par}
  \vspace{1.8cm}
  \begin{keybox}[项目定位]
  本文档包面向“谱图到分子结构”反问题的最小可执行科学工作流展示：以多模态谱图为输入，先完成可追踪的候选生成/检索与前向谱图一致性重排闭环，再为未来的图扩散、pairwise distance 与 TFN-Transformer 三维几何建模预留清晰扩展路径。
  \end{keybox}
  \vfill
  {\large 提交对象：DP Technology / Deep Potential Internship Application\par}
  \vspace{0.4cm}
  {\large 文档生成日期：\today\par}
\end{titlepage}

\pagenumbering{roman}
\section*{摘要}
\addcontentsline{toc}{section}{摘要}
实验谱图到分子结构的反演是化学分析中的核心难题，其本质是受噪声、模态缺失与一谱多解影响的逆问题。\projectname{} 将输入从单一谱图扩展为 IR/Raman/NMR/MS 等多模态证据，并将结构输出从单条 SMILES 提升为“候选分子图 + 前向一致性重排”的可验证工作流。当前 demo 重点验证四个环节：横坐标感知的多模态谱图编码、候选分子生成/检索、前向谱图一致性评分以及 Top-K 可视化与 trace log。未来完整版本将进一步引入可增删节点的 2D 图扩散、pairwise distance 矩阵预测、距离几何三维初始化，以及 parity-aware TFN-Transformer 精修模块，用以统一处理几何约束、张量特征与手性信息。本文档包的目标不是声称完整大模型已经训练完成，而是给出一条最小可执行、可追踪、可扩展的 AI for Science demo 路线。 

\begin{keybox}[摘要要点]
\begin{itemize}
  \item 问题定义：多模态谱图驱动的小分子结构反演；
  \item 当前目标：验证“编码--候选--前向一致性--Top-K”最小闭环；
  \item 未来扩展：graph diffusion、distance geometry、TFN-Transformer；
  \item Demo 价值：可执行、可审计、适合实习申请场景的科学工作流展示。
\end{itemize}
\end{keybox}

\tableofcontents
\clearpage
\pagenumbering{arabic}

\section{项目提议（500字内）}
\begin{keybox}
\textbf{具体目标}：构建一个面向 IR/Raman 的多模态谱图结构反演 demo，输入实验谱图，输出 Top-K 候选分子、谱图一致性分数与可视化结果，并保留完整 trace log。\par
\textbf{核心价值}：相较直接谱图到 SMILES，本方案将分子视为图与几何对象，用横坐标感知编码和前向谱图重排提升可解释性、可验证性与后续 3D 扩展能力。\par
\textbf{可行性简析}：先以公开小样本和候选库实现“编码--检索/轻量生成--一致性重排”最小闭环，再逐步扩展到图扩散、distance geometry 与 TFN 模块；数据、工具与评测路径明确，适合短周期 demo 验证。
\end{keybox}

\section{背景与动机}
实验谱图是分子结构解析中最重要的证据来源之一：IR/Raman 反映键振动与官能团环境，NMR 提供局部化学环境与连接关系线索，MS 则对分子式、碎片与分子量范围给出约束。近年的数据集与模型工作表明，多模态联合输入能够显著增强结构解析上限，而前向谱图预测模型已经具备为候选分子提供快速物理一致性评分的能力 \citep{Alberts2024MultimodalDataset,RocabertOriols2025VibraCLIP,Zou2023DetaNet}。

但将谱图直接翻译为 SMILES 仍存在三类关键局限。第一，SMILES 是人为序列，不是分子本体表示；同一分子可对应多种等价书写方式。第二，谱图到结构本质上是一谱多解的逆问题，单一序列输出会压缩不确定性表达。第三，谱图的物理来源依赖分子图、三维几何与张量响应；若忽略几何与前向验证，模型容易在实验域噪声下失去可解释性。以 IR/Raman 到结构的双向翻译工作为例，直接“谱图到结构”的范式虽有进展，但实验域性能、结构歧义与候选可验证性依然是主要瓶颈 \citep{Hu2025TranSpec,Alberts2024IRStructure}。

因此，本项目不把 demo 目标设定为“训练一个完整的大模型”，而是构造一条更符合科学反演逻辑的工作流：先从多模态谱图中提取结构证据，再生成或检索候选分子，最后通过前向谱图一致性与化学有效性完成重排。这样的设计更适合 DP Technology / Deep Potential 的 AI for Science 场景：它强调可追踪、可验证、可扩展，而不是只追求端到端生成结果的表面准确率。

\begin{keybox}[为什么适合作为实习 Demo]
\begin{itemize}
  \item 问题清晰：谱图到结构反演是典型的科学逆问题；
  \item 路径可落地：可先做最小闭环，再逐步引入更强生成与几何模块；
  \item 结果可验证：Top-K 输出、前向谱图一致性与 trace log 均可直接检查；
  \item 技术外延强：可自然扩展至图扩散、几何深度学习与物理先验融合。
\end{itemize}
\end{keybox}

\section{Demo 目标与边界}
本 demo 的目标不是直接声称“完整的 MultiSpec-GeoDiff / MultiSpec-GeoDiff-TFN 已训练完成”，而是验证一条最小可执行且具备科研解释力的谱图反演闭环。当前版本优先实现的内容如下：
\begin{enumerate}
  \item 支持 IR/Raman 为主的输入格式与基础预处理；
  \item 使用横坐标感知的谱图编码器生成全局条件向量与 token 级条件序列；
  \item 通过候选库检索或轻量候选生成器得到可重排的分子集合；
  \item 使用轻量前向谱图打分器与 RDKit 合法性规则完成 Top-K 重排；
  \item 生成可视化结果与 trace log，保证运行流程可复查。
\end{enumerate}

未来扩展部分则包括：可增删节点的 2D 图扩散、pairwise distance 预测、距离几何三维初始化、parity-aware TFN-Transformer 精修以及更强的 physics-informed forward spectrum prediction。文中所有这类内容都将明确标注为“未来扩展”或“完整版本路线”，而非当前已交付模块。

本次提交的预期输出包括：
\begin{itemize}
  \item 一套可编译的 LaTeX proposal package；
  \item 一段可直接粘贴到申请表中的 500 字内中文提议；
  \item 占位图、里程碑表与验证计划；
  \item 对后续代码 demo 的目录、输入输出与风险边界说明。
\end{itemize}

\begin{keybox}[不作出的声明]
\begin{itemize}
  \item 不宣称完整 diffusion 或 TFN 模型已经训练完成；
  \item 不宣称已经获得大规模实验结果或 SOTA 指标；
  \item 不把 Top-K 候选误写成唯一正确结构；
  \item 不把未来三维与手性模块描述成当前可用能力。
\end{itemize}
\end{keybox}

\section{方法总览}
\projectname{} 的核心思想是：不再把谱图看作需要直接翻译成分子字符串的“语言 token”，而是将其视为对分子图、几何与物理响应的多源观测。完整概念路径可写为
\[
\mathcal{S}_{\mathrm{multi}}
\rightarrow
(c, C)
\rightarrow
\mathcal{G}_{\mathrm{cand}}
\rightarrow
D
\rightarrow
R^{(0)}
\rightarrow
R
\rightarrow
\hat{\mathcal{S}}
\rightarrow
\mathrm{Top\mbox{-}K},
\]
其中 $\mathcal{S}_{\mathrm{multi}}$ 为多模态谱图，$(c, C)$ 分别表示全局条件向量与 token 级条件序列，$\mathcal{G}_{\mathrm{cand}}$ 为候选分子集合，$D$ 为 pairwise distance 矩阵，$R^{(0)}$ 为距离几何恢复的初始 3D 构象，$R$ 为等变精修后的几何表示，$\hat{\mathcal{S}}$ 为候选分子的前向预测谱图。

当前 demo 仅显式实现到 $\mathcal{G}_{\mathrm{cand}} \rightarrow \hat{\mathcal{S}} \rightarrow \mathrm{Top\mbox{-}K}$ 的最小闭环；但从文档设计上，后续 2D graph diffusion、pairwise distance 与 TFN-Transformer 的接口已经一并预留。这样既可以在短周期内提交一个可跑通的原型，也能展示项目向更完整科学模型演化的技术路线。

\begin{figure}[htbp]
  \centering
  \placeholderfigure{Figure placeholder: overall architecture}
  \caption{Overall architecture of MultiSpec-GeoDiff: multi-modal spectra are encoded, candidate molecules are generated or retrieved, forward spectra are predicted, and Top-K candidates are reranked.}
\end{figure}

\begin{keybox}[方法分层]
\begin{enumerate}
  \item \textbf{谱图层}：尊重真实横坐标与模态差异；
  \item \textbf{候选层}：先得到可验证的候选集合，而非单点猜测；
  \item \textbf{几何层}：用 pairwise distance 连接 2D 图与 3D 结构；
  \item \textbf{验证层}：通过前向谱图一致性与化学规则完成后验重排。
\end{enumerate}
\end{keybox}

\section{多模态谱图编码器}
编码器的设计目标是把不同仪器模态下的谱图统一为可供分子生成与重排使用的条件表示，同时尽量保留每个模态自身的物理坐标语义。对第 $m$ 个模态，记输入为
\[
S^{(m)} = \{(\nu_i^{(m)}, I_i^{(m)})\}_{i=1}^{L_m},
\]
其中 $\nu_i^{(m)}$ 是真实横坐标（如 cm$^{-1}$、ppm 或 $m/z$），$I_i^{(m)}$ 是对应强度。首先使用一维 CNN stem 提取局部峰形、肩峰、带宽与局部噪声模式，再叠加模态 embedding 与绝对横坐标编码：
\[
H_i^{(m,0)} = \operatorname{CNN}_m(S^{(m)})_i + e_m + p_m(\nu_i^{(m)}).
\]
这里的 $p_m(\nu)$ 由真实物理坐标而非 token 序号驱动，从而显式区分不同波数、化学位移或质荷比区域的结构意义。

模态内 attention 使用带相对谱图距离偏置的形式：
\[
A_{ij}^{(m)} = \operatorname{softmax}_j\left(\frac{Q_i^{(m)}K_j^{(m)\top}}{\sqrt{d}} + b_m\left(|\nu_i^{(m)}-\nu_j^{(m)}|\right)\right).
\]
该偏置反映同一模态内部峰间相对位置关系的重要性，例如 IR 中不同振动区间的组合证据往往比孤立峰值更稳定。跨模态 attention 则不直接引入坐标距离偏置：
\[
\widetilde{A}_{ij}^{(m\rightarrow n)} = \operatorname{softmax}_j\left(\frac{Q_i^{(m)}K_j^{(n)\top}}{\sqrt{d}}\right), \quad m \neq n,
\]
因为不同模态的横坐标处于不同物理坐标系中，直接比较绝对距离通常没有意义。跨模态层更适合学习“哪个模态的哪段证据可以互补解释另一个模态”的关系。

为了适应真实实验场景，编码器还需支持缺失模态输入。具体做法是：
\begin{itemize}
  \item 训练时使用 modality dropout；
  \item 推理时为缺失模态注入 mask token 与存在性标记；
  \item 输出同时包含全局条件向量 $c$ 与 token 级条件序列 $C$，便于后续候选生成与重排模块按需调用。
\end{itemize}

\begin{figure}[htbp]
  \centering
  \placeholderfigure{Figure placeholder: multi-modal spectral encoder}
  \caption{Coordinate-aware multi-modal spectral encoder. Intra-modal attention uses spectral-coordinate bias, while cross-modal attention does not use coordinate-distance bias.}
\end{figure}

\begin{keybox}[编码器输出]
\begin{itemize}
  \item \textbf{全局条件向量 $c$}：用于候选召回、检索与全局先验判断；
  \item \textbf{token 级条件序列 $C$}：用于候选细粒度打分或未来扩散过程中的逐步条件注入；
  \item \textbf{模态 mask 信息}：显式记录哪些证据缺失，避免模型错误地把“无数据”解释为“无信号”。
\end{itemize}
\end{keybox}

\section{候选生成与前向一致性重排}
MVP 阶段的关键任务不是直接训练完整分子扩散模型，而是尽快建立可验证的候选闭环。因此当前 demo 采用两类可替代策略：
\begin{enumerate}
  \item \textbf{候选库检索}：依据谱图编码向量、分子式约束、质量窗口或简单谱峰特征，从公开库中召回若干候选分子；
  \item \textbf{轻量候选生成}：在需要时引入规则驱动或小模型生成器补充候选多样性。
\end{enumerate}

候选集合生成后，先通过 RDKit 类化学规则执行基本过滤，包括价态合法性、芳香性/键类型合法性、分子式与质量范围匹配等，再交由前向谱图打分器进行一致性排序。对于 2D demo，前向打分器可以采用两种低风险实现：
\begin{itemize}
  \item 基于候选库已有谱图或近邻模板的相似度打分；
  \item 基于轻量 GNN 的近似前向谱图预测，再与 query 谱图计算余弦相似度或加权峰匹配得分。
\end{itemize}
重排输出不是未经验证的单个结构，而是带有解释分数的 Top-K 候选分子集合。这样更符合一谱多解的实际场景，也为后续人工审查或高保真计算留出了接口。

从系统演化角度看，当前的候选检索/轻量生成模块可视为未来可增删节点图扩散的受控代理：接口保持一致，但实现复杂度和计算成本显著更低。随着数据与算力条件成熟，后续可以把该模块替换为条件化的 insert/delete graph diffusion，而无需重写后续重排与验证层。

\begin{figure}[htbp]
  \centering
  \placeholderfigure{Figure placeholder: demo pipeline output}
  \caption{Minimum viable demo pipeline: query spectra, candidate library/generator, forward spectrum scorer, and Top-K visualization.}
\end{figure}

\begin{keybox}[重排评分建议]
可组合使用
\[
\mathrm{Score}(g)=\lambda_1\,s_{\mathrm{spec}}+\lambda_2\,s_{\mathrm{formula}}+\lambda_3\,s_{\mathrm{chem}}+\lambda_4\,s_{\mathrm{prior}},
\]
其中 $s_{\mathrm{spec}}$ 为谱图一致性，$s_{\mathrm{formula}}$ 为分子式/质量匹配，$s_{\mathrm{chem}}$ 为化学有效性，$s_{\mathrm{prior}}$ 为来自候选召回阶段的先验得分。该形式便于后续做模态权重学习与消融实验。
\end{keybox}

\section{未来扩展：几何扩散与 TFN-Transformer}
完整版本的 MultiSpec-GeoDiff-TFN 将在当前 demo 闭环之上补齐 2D 生成、2D-to-3D 桥接与等变几何精修三层能力。第一层是可增删节点的 2D 图扩散：相较固定节点数的离散生成器，这类方法更适合谱图条件下的未知分子大小推断 \citep{Ninniri2025GrIDDD}。第二层是 pairwise distance 矩阵预测，它不直接输出绝对坐标，而是先预测满足旋转/平移不变性的距离表示，再通过 MDS 或 distance geometry 恢复初始构象。第三层则是 parity-aware TFN-Transformer，用以在三维空间中联合更新坐标与张量特征。

从建模上看，pairwise distance 作为 2D 与 3D 之间的桥梁具有两个优势：一是它天然满足整体刚体变换不变性；二是它能与后续等变网络的相对距离、边偏置和局部几何编码自然衔接。初始构象恢复后，坐标更新可借鉴 EGNN 的相对坐标更新思想 \citep{Satorras2021EGNN}，而高阶特征表达与 attention 结构则可参考 Equiformer / EquiformerV2 一类基于 irreps 的等变 Transformer \citep{Liao2022Equiformer,Liao2024EquiformerV2}。

在 TFN-Transformer 模块中，可为每个原子维护 $0e,0o,1e,1o,2e,2o$ 等奇偶-阶数组合通道。直观地说，偶宇称标量/向量用于表达常规几何与化学环境，奇宇称 pseudoscalar 通道则为 CD/VCD 等手性敏感模态预留表达能力。attention 权重本身仍由不变量标量控制，但消息内容可由张量内积、边偏置与距离偏置共同调制。未来若接入更高保真的 forward model，则可进一步把 DetaNet 类张量谱预测能力并入后验打分环节 \citep{Zou2023DetaNet}。

需要强调的是，本节内容属于\emph{未来扩展}：它展示项目从可运行 demo 向完整科学模型延伸的合理路径，而非当前已完成实现。当前文档只要求把这些接口、科学动机与验证方向表达清楚。

\begin{figure}[htbp]
  \centering
  \placeholderfigure{Figure placeholder: future TFN-Transformer module}
  \caption{Future extension: insert/delete graph diffusion generates 2D graphs; pairwise distances initialize 3D geometry; TFN-Transformer refines coordinates and tensor features.}
\end{figure}

\begin{keybox}[未来完整路径]
{\small
\[
\begin{aligned}
&\text{multi-modal spectra}
\rightarrow
\text{insert/delete 2D graph diffusion}\\
&\rightarrow
\text{pairwise distance matrix}
\rightarrow
\text{distance geometry init}\\
&\rightarrow
\text{TFN-Transformer refinement}
\rightarrow
\text{physics-informed forward spectrum prediction}\\
&\rightarrow
\text{Top-K candidates}
\end{aligned}
\]
}
\end{keybox}

\section{Demo 实施计划与交付物}
文档层面，本次交付首先产出一个自包含的 XeLaTeX 文档包，用于固定 proposal 结构、图示位置、里程碑和验证标准；代码层面，则为后续 notebook / script demo 预留目录与输出接口。建议的仓库交付物分为两层：
\begin{itemize}
  \item \textbf{当前已生成的文档包}
  \begin{itemize}
    \item \texttt{latex/}
    \item 根目录 \texttt{README.md}
    \item \texttt{docs/Demo\_Documentation\_Build\_Report.md}
  \end{itemize}
  \item \textbf{后续 demo 运行产物}
  \begin{itemize}
    \item \texttt{outputs/topk\_candidates.csv}
    \item \texttt{outputs/spectrum\_match.png}
    \item \texttt{outputs/molecule\_grid.png}
    \item \texttt{outputs/trace\_log.json}
  \end{itemize}
\end{itemize}

\begin{keybox}[建议目录视图]
\begin{itemize}
  \item \texttt{latex/main.tex}
  \item \texttt{latex/sections/}
  \item \texttt{latex/figures/}
  \item \texttt{latex/tables/}
  \item \texttt{latex/refs.bib}
  \item \texttt{latex/Makefile}
\end{itemize}
\end{keybox}

建议将开发过程按里程碑推进：先稳定文档与接口，再补齐最小数据流、编码器 forward、候选生成与重排，可视化与 trace log 最后闭环。对应里程碑见表~\ref{tab:milestones}。

\begin{table}[htbp]
  \centering
  \caption{Demo 里程碑规划}
  \label{tab:milestones}
  \small
  \renewcommand{\arraystretch}{1.2}
  \begin{tabular}{>{\raggedright\arraybackslash}p{0.10\linewidth} >{\raggedright\arraybackslash}p{0.22\linewidth} >{\raggedright\arraybackslash}p{0.28\linewidth} >{\raggedright\arraybackslash}p{0.28\linewidth}}
    \toprule
    ID & 里程碑 & 目标 & 预期交付物 \\
    \midrule
    M1 & Document package & 固定 proposal 结构、图注、风险边界与构建方式 & LaTeX 文档包、README、build report \\
    M2 & Minimal data pipeline & 统一输入谱图格式、样例数据与预处理接口 & 输入样例、预处理脚本或 notebook 单元 \\
    M3 & Spectrum encoder forward pass & 跑通 coordinate-aware encoder 的前向流程 & 编码张量、trace log、模块说明 \\
    M4 & Candidate retrieval / generation & 建立候选召回与基础化学过滤闭环 & 候选列表、合法性过滤结果 \\
    M5 & Forward spectrum reranking & 使用近似前向谱图或模板完成 Top-K 重排 & Top-K csv、相似度分数、对比图 \\
    M6 & Visualization and trace log & 输出结果面板与可追踪日志 & 分子网格图、谱图叠图、trace\_log.json \\
    M7 & Optional graph diffusion prototype & 预留未来 insert/delete graph diffusion 原型接口 & 原型说明、伪代码或最小实验记录 \\
    \bottomrule
  \end{tabular}
\end{table}


此外，文档中所有图均先以占位框表示，方便后续替换为正式架构图、流程图和结果图。建议在实际补图时统一采用英文图内标签、中文正文解释的形式，保证技术感与国际化可读性兼顾。

\begin{figure}[htbp]
  \centering
  \placeholderfigure{Figure placeholder: expected demo outputs}
  \caption{Expected demo outputs: input spectrum, Top-K molecular candidates, spectrum overlay, and trace log.}
\end{figure}

\section{验证方案与主要风险}
验证部分分为\textbf{demo 闭环验证}与\textbf{未来模型验证}两层。当前 demo 阶段优先回答三个问题：第一，候选检索/生成是否能覆盖合理候选；第二，前向谱图一致性重排是否能提升候选排序质量；第三，整条流程是否留下足够完备的 trace 信息用于复查。围绕这三个问题，建议使用以下指标：
\begin{itemize}
  \item Top-K accuracy / recall；
  \item 谱图相似度（余弦相似度、峰匹配得分或相关系数）；
  \item 分子式/质量匹配率；
  \item 官能团匹配率；
  \item RDKit validity 与化学规则通过率；
  \item reranking gain（重排前后 Top-1 / Top-K 提升）。
\end{itemize}
对于未来完整模型，还可增加 pairwise distance MAE、3D RMSD、chirality consistency、构象覆盖率等几何指标。

主要风险包括：
\begin{enumerate}
  \item \textbf{实验噪声与域偏移}：实验谱图与模拟/公开小样本之间可能存在系统差异；
  \item \textbf{一谱多解歧义}：不同结构可能共享相似谱图特征；
  \item \textbf{数据稀缺}：多模态联合样本、尤其带高质量 3D/手性标注的数据有限；
  \item \textbf{过度承诺风险}：若不区分 demo 与 future module，容易在文档中形成不实暗示；
  \item \textbf{3D/手性信息不完整}：仅凭 2D 候选或单模态信息往往不足以完全恢复立体信息。
\end{enumerate}

针对上述风险，本文档建议采用以下缓解策略：
\begin{itemize}
  \item 以 Top-K 候选而非单解输出承载不确定性；
  \item 用前向谱图一致性和 trace log 提高结果可解释性；
  \item 在公开小规模可控数据上先做消融与闭环验证；
  \item 对所有“未来扩展”模块明确标注，避免把路线图写成已得结果；
  \item 在后续阶段引入模态 dropout、候选覆盖分析与人工审查接口。
\end{itemize}

\begin{keybox}[推荐验证顺序]
\begin{enumerate}
  \item 先验证输入--候选--重排闭环是否可运行；
  \item 再比较“无重排 / 有重排”的排序提升；
  \item 最后在小样本上分析失败案例与缺模态鲁棒性。
\end{enumerate}
\end{keybox}

\section{结论}
\projectname{} 适合作为 DP Technology / Deep Potential 实习申请中的 demo proposal，原因在于它同时满足四个条件：第一，问题具有明确的科学意义，直接对应实验谱图到分子结构反演这一高价值逆问题；第二，方案具备最小可执行路径，可以先交付一个编码--候选--前向一致性重排的闭环原型；第三，方法论与 AI for Science 高度契合，强调图结构、几何约束、物理一致性与可追踪验证；第四，体系具有自然扩展性，后续可以平滑过渡到 graph diffusion、pairwise distance 与 TFN-Transformer 等更强的几何深度学习模块。

从实习 demo 的角度看，这种“先做可运行闭环，再逐步补强科学深度”的策略比直接承诺一个完整大模型更可信，也更能体现对问题本质和工程边界的理解。若后续补齐真实图示、最小数据流和样例输出，本提案即可作为一份兼具研究视角、技术路线和执行可行性的申请材料。


\clearpage
\bibliographystyle{unsrtnat}
\bibliography{refs}

\end{document}
